\begin{theorem}[导子 \(2\) 的可见性、导子 \(4\) 的 Ramanujan 失明与 primitive quarter-turn 的选相]\label{thm:conclusion-window6-conductor-two-quarterturn-selection}
\leanverified{paper\_conclusion\_window6\_conductor\_two\_quarterturn\_selection}
对每个 \(w\in X_6\)，令 \(\Delta_w\) 为 centered hidden template \(D_w\) 上的均匀随机变量。则有：
\[
\mathbb E\,c_2(\Delta_w)=\frac{4}{d(w)}-1,
\]
亦即
\[
\mathbb E\,c_2(\Delta_w)=
\begin{cases}
1,& d(w)=2,\\[2mm]
\frac13,& d(w)=3,\\[2mm]
0,& d(w)=4;
\end{cases}
\]
\[
\mathbb E\,c_4(\Delta_w)=0
\qquad (\forall\, w\in X_6);
\]
若记 primitive conductor-\(4\) 角色
\[
\chi_4(n):=i^n,
\]
则
\[
\mathbb E\,\chi_4(\Delta_w)=
\begin{cases}
0,& d(w)=2,\\[2mm]
\frac{i}{3},& d(w)=3,\\[2mm]
0,& d(w)=4.
\end{cases}
\]
从而
\[
\bigl(\mathbb E\,c_2(\Delta_w),\,\mathbb E\,\chi_4(\Delta_w)\bigr)
\in
\Bigl\{(1,0),\ \bigl(\tfrac13,\tfrac{i}{3}\bigr),\ (0,0)\Bigr\},
\]
并且这三值恰与
\[
d(w)=2,3,4
\]
一一对应。换言之，完整的 \(c_4\) 对局部 hidden template 全部失明，而 primitive quarter-turn 切片只在三点模板上存活；与导子 \(2\) 通道联用时，最低两个 \(2\)-adic 角色已足以完全分辨三类局部 hidden 相。
\end{theorem}
\begin{proof}
由定义，
\[
c_2(n)=\sum_{\substack{a\bmod 2\\(a,2)=1}}e^{2\pi i an/2}=(-1)^n.
\]
对三类模板分别计算可得
\[
\frac{1+1}{2}=1,\qquad
\frac{1-1+1}{3}=\frac13,\qquad
\frac{1-1+1-1}{4}=0,
\]
故第一组结论成立。

再由
\[
c_4(n)=\sum_{\substack{a\bmod 4\\(a,4)=1}}e^{2\pi i an/4}=i^n+(-i)^n,
\]
对
\[
D_w\in\bigl\{\{0,34\},\ \{0,21,34\},\ \{0,21,34,55\}\bigr\}
\]
逐一求平均，得到
\[
\frac{2+(-2)}{2}=0,\qquad
\frac{2+0-2}{3}=0,\qquad
\frac{2+0-2+0}{4}=0.
\]
故
\[
\mathbb E\,c_4(\Delta_w)=0
\qquad (\forall\, w\in X_6).
\]

最后，对 primitive 角色
\[
\chi_4(n)=i^n
\]
分别计算
\[
\frac{1+(-1)}{2}=0,\qquad
\frac{1+i-1}{3}=\frac{i}{3},\qquad
\frac{1+i-1-i}{4}=0.
\]
三值彼此不同，因此
\[
\bigl(\mathbb E\,c_2(\Delta_w),\,\mathbb E\,\chi_4(\Delta_w)\bigr)
\]
与三种局部模板类型一一对应。证毕。
\end{proof}

\begin{theorem}[有限体积下导子 \(2\) 与 primitive 导子 \(4\) 的精确相选择律]\label{thm:conclusion-window6-finitevolume-conductor-two-quarterturn-selection}
设 \(\Lambda\) 为有限指标集，\(w=(w_x)_{x\in\Lambda}\in X_6^\Lambda\)。在 \(\pi_\Lambda^{-1}(w)\) 上取均匀测度，并定义总 centered hidden offset
\[
\Delta_\Lambda:=\sum_{x\in\Lambda}\Delta_{w_x}.
\]
再记
\[
N_j(w):=\#\{x\in\Lambda:d(w_x)=j\}
\qquad (j=2,3,4).
\]
则有精确因子化公式
\[
\mathbb E_w\,c_2(\Delta_\Lambda)
=
3^{-N_3(w)}\,\mathbf 1_{\{N_4(w)=0\}},
\]
以及
\[
\mathbb E_w\,\chi_4(\Delta_\Lambda)
=
\left(\frac{i}{3}\right)^{N_3(w)}\mathbf 1_{\{N_2(w)=N_4(w)=0\}}.
\]
因此，只要出现一个四点纤维位点，整个全局奇偶 Ramanujan 通道便严格湮灭；而 primitive conductor-\(4\) 通道存活当且仅当每一个位点都处于三点模板，它因而是 window-\(6\) hidden sector 的纯三点相精确选择器。
\end{theorem}
\begin{proof}
由有限体积分解，
\[
\pi_\Lambda^{-1}(w)
=
\prod_{x\in\Lambda}\bigl(V_6(w_x)+D_{w_x}\bigr),
\]
故在均匀测度下，各坐标 centered offset \(\Delta_{w_x}\) 相互独立。又由于
\[
c_2(n)=(-1)^n,\qquad \chi_4(n)=i^n
\]
都是加法角色，遂有
\[
c_2\!\left(\sum_{x\in\Lambda}\Delta_{w_x}\right)=\prod_{x\in\Lambda}c_2(\Delta_{w_x}),
\qquad
\chi_4\!\left(\sum_{x\in\Lambda}\Delta_{w_x}\right)=\prod_{x\in\Lambda}\chi_4(\Delta_{w_x}).
\]
由独立性得
\[
\mathbb E_w\,c_2(\Delta_\Lambda)=\prod_{x\in\Lambda}\mathbb E\,c_2(\Delta_{w_x}),
\qquad
\mathbb E_w\,\chi_4(\Delta_\Lambda)=\prod_{x\in\Lambda}\mathbb E\,\chi_4(\Delta_{w_x}).
\]
再代入定理 \ref{thm:conclusion-window6-conductor-two-quarterturn-selection} 中的局部三值表
\[
\mathbb E\,c_2(\Delta_{w_x})=
\begin{cases}
1,& d(w_x)=2,\\
\frac13,& d(w_x)=3,\\
0,& d(w_x)=4,
\end{cases}
\qquad
\mathbb E\,\chi_4(\Delta_{w_x})=
\begin{cases}
0,& d(w_x)=2,\\
\frac{i}{3},& d(w_x)=3,\\
0,& d(w_x)=4,
\end{cases}
\]
便立即得到两条精确公式。证毕。
\end{proof}

\begin{proposition}[所有良性奇导子都严格非消失]\label{prop:conclusion-window6-all-good-odd-conductors-strictly-positive}
\leanverified{paper\_conclusion\_window6\_all\_good\_odd\_conductors\_strictly\_positive}
若 \(q>1\) 为奇整数且
\[
(q,39270)=1,
\]
则对任意 \(w\in X_6\) 都有
\[
\mathcal R_w(q)>0.
\]
特别地，在与 \(39270\) 互素的导子范围内，任何奇导子 Ramanujan 通道都不可能产生局部精确消失；window-\(6\) hidden sector 的精确 annihilation 因而纯属 \(2\)-adic 现象。
\end{proposition}
\begin{proof}
由
\[
(q,39270)=1
\]
且 \(q\) 为奇整数可知，\(q\) 的最小可能值为 \(13\)，故
\[
\varphi(q)\ge \varphi(13)=12.
\]
另一方面，由定理 \ref{thm:conclusion-window6-hidden-template-primegood-ramanujan-faultlaw}，
\[
d(w)\mathcal R_w(q)=\varphi(q)+(d(w)-1)\mu(q).
\]
由于
\[
d(w)-1\le 3,\qquad \mu(q)\in\{-1,0,1\},
\]
遂有
\[
d(w)\mathcal R_w(q)\ge \varphi(q)-3\ge 9>0.
\]
故
\[
\mathcal R_w(q)>0.
\]
证毕。
\end{proof}

\begin{conjecture}[prime-to-\(39270\) 的 Ramanujan 通道只编码退化度场，全部精确缺口由 \(2\)-adic 例外穷尽]\label{conj:conclusion-window6-primegood-ramanujan-channels-only-degeneracy}
设一族 canonical sitewise window-\(6\) lifts 处于可见 sector 的高斯自由场极限之下。则对任意固定导子 \(q\) 满足
\[
(q,39270)=1,
\]
在去除 visible 平移相位之后，hidden sector 的 Ramanujan 观测在连续极限中只能收敛到局部退化度场
\[
d_x:=d(w_x)\in\{2,3,4\}
\]
的函数；尤其不能创造新的红外零模，也不能生成新的连续精确缺口。全部精确频域缺口应由两类且仅两类 \(2\)-adic 例外穷尽：导子 \(2\) 的奇偶通道与 primitive 导子 \(4\) 的 quarter-turn 通道。
\end{conjecture}
